% Ad Familiares 5.19 --- headnote
% ad-familiares-05-19, around 28 April 49 BC, Cumae

\textit{Cicero to M. Mescinius Rufus, written from the villa at
Cumae around 28 April 49 BC --- Perseus dateline \textit{Scr.
in Cumano circ. iv K. Mai. a. 705 (49)}. The salutation
\textit{CICERO RVFO} is the same as that of Fam. 5.20 and of
the earlier letters in this block; the recipient is Cicero's
quaestor of 51--50, with whom he has just finished settling
provincial accounts (Fam. 5.20), and not (as the manuscript
heading sometimes suggests) Manius Curius. Cicero is now in
Campania, three weeks before he will sail to join Pompey, and
the question consuming him is exactly the question Sulpicius
raised in Fam. 4.1: whether and how to leave Italy.}

\textit{Mescinius has written to say that, free now of any
obligation owed to a governor in office, he is more devoted to
Cicero than ever, and ready to follow him in whatever course he
takes. Cicero seizes the offer with a Stoic frame: ``what is
right is plain; what is expedient is dark, though even so, if
we are worthy of our studies and our writings, we cannot doubt
that those courses are most useful which are most right'' ---
\textit{quid rectum sit apparet, quid expediat obscurum est}.
He accepts the offer, but with care: if you can't come, fear
made you say one thing and your decency the other, and I will
forgive you. The decision he refers to (taken ``now for some
time'') is to join Pompey in Epirus, which he will do by
mid-June. The letter is the most directly Stoic of the April
correspondence; Cicero is rehearsing publicly, to a friend, the
calculation he has already made privately.}
